\usetikzlibrary{calc}
\begin{figure}[tb]
\definecolor{bitcolor}{gray}{0.80}
\newlength{\bitsize}
\setlength{\bitsize}{1.0cm}
\newlength{\meaningheight}
\setlength{\meaningheight}{-1.5cm}
\newlength{\xpos}
\setlength{\xpos}{10.6cm}
\newlength{\ypos}
\setlength{\ypos}{-1.5cm}
\begin{tikzpicture}
\small\sffamily
%
% First the bits
\begin{scope}[minimum size=\bitsize,very thick]
\node (b11) at (0\bitsize,0pt) [rectangle, draw, fill=bitcolor, label=above:10] {\small\ttfamily 1024};
\node (b9) at (1\bitsize,0pt) [rectangle, draw, fill=bitcolor, label=above:9] {\small\ttfamily 512};
\node (b8) at (2\bitsize,0pt) [rectangle, draw, fill=bitcolor, label=above:8] {\small\ttfamily 256};
\node (b7) at (3\bitsize,0pt) [rectangle, draw, fill=bitcolor, label=above:7] {\small\ttfamily 128};
\node (b6) at (4\bitsize,0pt) [rectangle, draw, fill=bitcolor, label=above:6] {\small\ttfamily 64};
\node (b5) at (5\bitsize,0pt) [rectangle, draw, fill=bitcolor, label=above:5] {\small\ttfamily 32};
\node (b4) at (6\bitsize,0pt) [rectangle, draw, fill=bitcolor, label=above:4] {\small\ttfamily 16};
\node (b3) at (7\bitsize,0pt) [rectangle, draw, fill=bitcolor, label=above:3] {\small\ttfamily 8};
\node (b2) at (8\bitsize,0pt) [rectangle, draw, fill=bitcolor, label=above:2] {\small\ttfamily 4};
\node (b1) at (9\bitsize,0pt) [rectangle, draw, fill=bitcolor, label=above:1] {\small\ttfamily 2};
\node (b0) at (10\bitsize,0pt) [rectangle, draw, fill=bitcolor, label=above:0] {\small\ttfamily 1};
% A key
\node at (11\bitsize,0pt) [rectangle, minimum size=\bitsize, label=above:Bit, align=left] {\small\ttfamily\null\hspace{1em} Value};
\end{scope}
%
% Now the meanings
\newcommand{\meaningbox}[2]{\parbox{5cm}{\raggedright\textbf{#1} #2}}
\begin{scope}[inner sep=0pt, outer sep=4pt, anchor=west]
\node (m0) at (\xpos,\ypos+0\meaningheight) {\meaningbox{Flag -- Bad Profile:}%
  {Do not use this profile (see ``information'' bits for an explanation).}};
\node (m1) at (\xpos,\ypos+1\meaningheight) {\meaningbox{Flag -- Warning:}%
  {This profile is questionable (see ``information'' bits for an explanation).}};
\node (m2) at (\xpos,\ypos+2\meaningheight) {\meaningbox{Flag -- Comment}%
  {See ``information'' bits for additional comments concerning this profile.}};
%
% Nudge a bit
%
\addtolength{\ypos}{-2mm}
\node (m4) at (\xpos,\ypos+3\meaningheight) {\meaningbox{Information:}%
  {(Warning) This profile may have been affected by high altitude clouds.}};
\node (m5) at (\xpos,\ypos+4\meaningheight) {\meaningbox{Information:}%
  {(Warning) This profile may have been affected by low altitude clouds.}};
\node (m6) at (\xpos,\ypos+5\meaningheight) {\meaningbox{Information:}%
  {(Comment) GEOS-5 \emph{a priori} temperature data were unavailable for this profile.}};
\node (m7) at (\xpos,\ypos+6\meaningheight) {\meaningbox{Information:}%
  {(Bad profile) The retrieval for this phase encountered a numerical error.}};
\node (m8) at (\xpos,\ypos+7\meaningheight) {\meaningbox{Information:}%
  {(Bad profile) Too few radiances were available for good retrieval of this profile.}};
\node (m9) at (\xpos,\ypos+8\meaningheight) {\meaningbox{Information:}%
  {(Bad profile) The task retrieving this profile crashed (typically a computer failure).}};
\end{scope}
%
% Now a separator
\draw[ultra thick] ($ (m2.west)!0.5!(m4.west) $) -- ($ (m2.east)!0.5!(m4.east) $);
% \draw[ultra thick] (0,0) -- (10,0);
%
% Now the connections
\begin{scope}[>=stealth,thick]
\draw[->] (b0) |- (m0);
\draw[->,dashed] (b1) |- (m1);
\draw[->] (b2) |- (m2);
\draw[->,dashed] (b4) |- (m4);
\draw[->] (b5) |- (m5);
\draw[->,dashed] (b6) |- (m6);
\draw[->] (b7) |- (m7);
\draw[->,dashed] (b8) |- (m8);
\draw[->] (b9) |- (m9);
\end{scope}
\end{tikzpicture}
\caption{The meaning of the various bits in the \texttt{Status} field.
  The bits not labeled are not used in \thisv.  Later versions may
  implement specific meanings for these bits.  Note that bit 6 (GEOS-5
  data) was not used in v1.5, and that the information in bits 7 and 8
  were combined into bit 8 in versions~1.5 and~2.2.}
\label{fig:StatusBits}
\end{figure}


%%% Local Variables: 
%%% mode: latex
%%% TeX-master: "v4-x-report"
%%% End: 
